\section{Introdução}
\label{sec:introducao}

As plataformas de empréstimos \textit{peer-to-peer} põem em contacto direto quem pede dinheiro e quem o investe. A LendingClub, a maior plataforma norte-americana deste tipo, decide para cada pedido se o aceita e que grau de risco e taxa de juro lhe atribui; quando o mutuário deixa de pagar, são os investidores que suportam a perda. Prever bem a taxa e o risco de incumprimento é, por isso, central tanto para a plataforma como para quem investe.

Este relatório corresponde à Parte~1 do projeto e cobre as três primeiras fases do ciclo de vida da análise de dados: formulação do problema, análise e limpeza dos dados e seleção de modelos. Usa-se o conjunto público \textit{Lending Club Loan Data (2007--2018)}~\cite{kaggle}, com 2\,260\,668 empréstimos financiados e 27\,648\,741 pedidos recusados. Definem-se três tarefas: uma regressão, que prevê a taxa de juro que seria atribuída a um pedido antes de ser fixada; uma classificação, que prevê, no momento da concessão, se o empréstimo vai ser pago ou entrar em \textit{default}; e um clustering de mutuários e de estados. Os modelos são treinados e avaliados na Parte~2.

O principal risco num problema deste tipo é a fuga de informação (\textit{data leakage}): o ficheiro contém o historial completo de cada empréstimo, incluindo os pagamentos posteriores à concessão, e um modelo que use essa informação parece bom na avaliação, mas é inútil na prática. Por isso, o trabalho segue quatro princípios, fixados antes de olhar para os dados: cada variável só é usada se já existir no momento de previsão da tarefa; a divisão em treino e teste é temporal, simulando um modelo treinado no início de 2015 e aplicado aos empréstimos desse ano; o teste final só é usado na avaliação final; e todas as transformações aprendidas com os dados são ajustadas apenas no treino, dentro de \textit{pipelines}. O código está organizado num pacote Python orientado a objetos, usado pelos \textit{notebooks}, com sementes fixas, para que os resultados sejam reprodutíveis. Usaram-se o scikit-learn~\cite{sklearn} nos \textit{pipelines} e o PCA e o UMAP~\cite{umap} na redução de dimensionalidade; na Parte~2 serão usados, entre outros, o LightGBM~\cite{lightgbm} e redes neuronais em Keras/TensorFlow~\cite{geron}.

O relatório está organizado da seguinte forma. A Secção~\ref{sec:formulacao} formula o problema: tarefas, critérios de sucesso, momentos de previsão, elegibilidade das variáveis, coortes e divisão temporal. A Secção~\ref{sec:analise} descreve o pré-processamento, a análise exploratória e os testes de hipóteses. A Secção~\ref{sec:modelos} apresenta as \textit{features} criadas, os modelos candidatos e o plano de validação, de métricas e de orçamento computacional. A Secção~\ref{sec:conclusao} resume as conclusões, as limitações e os próximos passos.